% Generated draft: requires mathematical and source review.
\documentclass[11pt]{article}
\usepackage[margin=25mm]{geometry}
\usepackage{fontspec}
\setmainfont{Latin Modern Roman}
\usepackage{amsmath,amssymb,mathtools,unicode-math}
\setmathfont{Latin Modern Math}
\usepackage{enumitem,needspace,xurl,hyperref}
\setlength{\parindent}{0pt}
\setlength{\parskip}{5pt}
\setlength{\emergencystretch}{3em}
\newcommand{\N}{\mathbb{N}}
\newcommand{\Z}{\mathbb{Z}}
\newcommand{\Q}{\mathbb{Q}}
\newcommand{\R}{\mathbb{R}}
\newcommand{\C}{\mathbb{C}}
\newcommand{\F}{\mathbb{F}}
\newcommand{\A}{\mathbb{A}}
\newcommand{\PP}{\mathbb{P}}
\newcommand{\E}{\mathbb{E}}
\newcommand{\eps}{\varepsilon}
\newcommand{\dd}{\,\mathrm{d}}
\begin{document}
\section*{UnsolvedMath ID: 701186}
\textbf{Source problem code:} catalogue entry
\subsection{Definitions and mathematical statement}
A formation \(\mathfrak F\) is a nonempty class of finite groups, closed under isomorphisms, homomorphic images, and finite subdirect products. It is saturated if, for every finite group \(H\), membership \(H/\Phi(H)\in\mathfrak F\) implies \(H\in\mathfrak F\), where \(\Phi(H)\) is the intersection of the maximal subgroups of \(H\). Define the \(\mathfrak F\)-residual by
\[
H^{\mathfrak F}=\bigcap\{N\mathrel{\unlhd}H:H/N\in\mathfrak F\}.
\]
A finite group \(G\) is minimal non-\(\mathfrak F\) if \(G\notin\mathfrak F\) and every proper subgroup of \(G\) belongs to \(\mathfrak F\). A group is soluble if its derived series terminates at the identity subgroup. A Sylow \(p\)-subgroup of a finite group has order equal to the largest power of the prime \(p\) dividing the group order.

Let \(\mathfrak F\) be any saturated formation, let \(p\) be a prime, and let \(G\) be any finite soluble minimal non-\(\mathfrak F\) group such that \(G^{\mathfrak F}\) is a Sylow \(p\)-subgroup of \(G\). The source asks whether \(G^{\mathfrak F}\) must be isomorphic to a factor group of a Sylow \(p\)-subgroup of the group it denotes by \(\operatorname{PSU}(3,p^{2n})\).

Under the intended but unstated existential interpretation of \(n\), the question is whether there always exist an integer \(n\geq 1\), a Sylow \(p\)-subgroup \(S\) of \(\operatorname{PSU}(3,p^{2n})\), and a normal subgroup \(K\unlhd S\) such that
\[
G^{\mathfrak F}\cong S/K.
\]
Here \(n\), \(S\), and \(K\) may depend on the given formation and group. This interpretation of the quantifier on \(n\) requires confirmation: the supplied statement neither quantifies \(n\) nor specifies its range.

The notation \(\operatorname{PSU}(3,p^{2n})\) is retained exactly in its source parameterization. It denotes a projective special unitary group, but the source record does not specify whether \(p^{2n}\) is the field size on which the unitary matrices are defined or the unitary parameter whose square is that field size. This convention must be resolved before replacing the notation with a fully specified group construction.

The supplied source asserts that the answer is affirmative when \(\mathfrak F\) is the formation of all finite nilpotent groups, and attributes that special case to the two references below. This assertion is retained as supplied background, without independent verification.
\subsection{Short English statement}
For a saturated formation \(\mathfrak F\), consider a finite soluble group outside \(\mathfrak F\) whose proper subgroups all lie in \(\mathfrak F\). If its \(\mathfrak F\)-residual is a Sylow \(p\)-subgroup, must that residual be a quotient of a Sylow \(p\)-subgroup of \(\operatorname{PSU}(3,p^{2n})\) for some positive integer \(n\)? The quantifier on \(n\) and the unitary-group parameter convention require confirmation. The source reports an affirmative answer for the formation of finite nilpotent groups.
\subsection{Sources}
\begin{itemize}
\item Local UnsolvedMath registry, permanent ID \(701186\), exact source problem code \texttt{catalogue entry}, title \emph{Kourovka Notebook Problem 15.98}. The supplied record labels the problem open; that label is not an independently verified literature assessment.
\item L. A. Shemetkov, proposer, \emph{The Kourovka Notebook (Unsolved Problems in Group Theory)}, Problem 15.98, 15th issue (2002). The supplied background also identifies the 21st issue edition (2026), arXiv:1401.0300v48. The supplied collection description names E. I. Khukhro and V. D. Mazurov as editors. These source details have not been independently checked.
\item V. D. Mazurov and S. A. Syskin, \emph{Math. Notes} \textbf{14}, no. 2 (1973), 683--686. Supplied as a reference for the finite nilpotent formation special case; no article title or link is supplied.
\item A. Kh. Zhurtov and S. A. Syskin, \emph{Siberian Math. J.} \textbf{26}, no. 2 (1987), 235--239. Bibliographic details are reproduced as supplied. Supplied as a reference for the finite nilpotent formation special case; no article title or link is supplied.
\end{itemize}
\end{document}
